\subsection{The seven-stage construction and its pulse convention}
\label{sv:reciprocal-template}
The reciprocal theorem in Section~\ref{sv:imports} uses a source
selector to assign each rectangle its own height, performs its planar
map there, and uses a target selector to return it to the coding height.
We write the fields from its proof~\cite[Theorem~3.2]{Entry} in a form
allowing separate shear and translation selectors, since later
constructions change individual profiles or stages.

Write the centers of $P_i,Q_i$ as $c_i^-,c_i^+$ and the prescribed
linear part as $\operatorname{diag}(r_i,r_i^{-1})$.  Let $z_0$ be
the coding height.  Choose smooth periodic selectors $A_i^-,A_i^+$
with zero planar mean, equal to one near their own source or target
rectangle and zero near all other rectangles of the same family.
Choose distinct private heights $z_i$.  Two height selectors
$\sigma_i,\tau_i$ equal one near $z_i$ and zero near every other
private height.  Require $\tau_i$ to have zero mean.  For $j=1,2$,
choose a periodic profile $g_{ij}$ equal to $s-c_{ij}^-$ on a
neighborhood of the coordinate range of the four centered shear paths,
and require
\[
 \left(\int\sigma_i\right)\left(\int g_{ij}\right)=0.
\]
Thus the shear selector $\sigma_i$ may have nonzero mean when both
linear profiles have zero mean.  Usually one can take
$\sigma_i=\tau_i$.  The derivative, odd-profile, and translated-copy
recipes above provide these choices; their compensating parts must
miss the other tracked rectangles or private-height plateaus as
specified there.

For the chronological shear list~\eqref{sv:list-a}, the seven spatial
fields are
\begin{align}
 W_1&=e_z\sum_i(z_i-z_0)A_i^-(x,y),\notag\\
 W_2&=e_y\sum_i(-r_i)\sigma_i(z)g_{i1}(x),&
 W_3&=e_x\sum_i(r_i^{-1}-1)\sigma_i(z)g_{i2}(y),\notag\\
 W_4&=e_y\sum_i\sigma_i(z)g_{i1}(x),&
 W_5&=e_x\sum_i(r_i-1)\sigma_i(z)g_{i2}(y),\notag\\
 W_6&=\sum_i\tau_i(z)(c_i^+-c_i^-,0),\notag\\
 W_7&=e_z\sum_i(z_0-z_i)A_i^+(x,y).
 \label{sv:seven-fields}
\end{align}
At each stage all branch contributions have the same direction of
motion, except in $W_6$, whose two horizontal components depend only
on height.  Consequently each entire stage is transverse, even where
profile supports overlap.  The stated integral conditions give zero
spatial mean to every summand and hence to each $W_j$.

To turn the stages into smooth motion, choose successive closed
rational intervals $[a_j,b_j]$ with $0<a_j<b_j<1$ and positive gaps,
and set
\begin{equation}\label{sv:pulse-convention}
 \beta_j(t)=\frac1{b_j-a_j}
 H'\!\left(\frac{t-a_j}{b_j-a_j}\right),\qquad
 U(t,x)=\sum_j\beta_j(t)W_j(x)\quad(0\le t\le1).
\end{equation}
The step function $H$ defined above is nondecreasing, so each pulse
is nonnegative, has integral one, and vanishes to every order at its
endpoints.  The finite list repeats smoothly with period one and zero
collars at integer times.  This convention applies to any finite number
of repeated stages below.  The same pulse formula is used on a
separately specified finite loading interval.  Its derivatives and
their bounds are effective by the same flat-profile estimates.  Distinct stages have disjoint temporal
supports; branch contributions within one stage act simultaneously.
Thus the summed velocity still has zero self-advection, and
Lemma~\ref{sv:force} applies.

Here is the exact path represented by these fields.  Starting from
$(y,z_0)$ with $y\in P_i$, $W_1$ lifts the point to $(y,z_i)$.
Only the $i$th height selectors act during the next five stages.
The four shears send the offset $y-c_i^-$ to
$\operatorname{diag}(r_i,r_i^{-1})(y-c_i^-)$, still based at
$c_i^-$.  The sixth stage translates that center to $c_i^+$, and
$W_7$ lowers the resulting point of $Q_i$ back to $z_0$.
Nonnegative pulse masses parameterize each partial shear or translation
by a number in $[0,1]$.  The excursion estimates above therefore
verify the linear plateaus on the full closed source rectangle before
ODE uniqueness identifies these formulas with the actual trajectories.
Positive plateau margins also permit a sufficiently small neighborhood
of that rectangle and coding height; its size must be chosen from the
margins and the affine partial maps.  The bounds for points on the
original rectangle are not automatically bounds for that neighborhood.

The other four-shear lists replace $W_2,\ldots,W_5$ by their listed
coefficients and directions, using their own excursion estimates.
Splitting $W_6$ into an $x$-translation and a $y$-translation gives the
eight-stage version.  If extra transverse controls are used, the first
is one on the current $y$-interval and the second on the target
$x$-interval already reached; both retain the private-height selector.
Their scalar coefficients must have zero integral, supplied by either
factor or by a correction supported off every tracked path.  Remote
centers and routes with an extra translation are specified separately
below: their endpoint maps use the same principles, while their
intermediate paths require the bounds proved for those routes.
